\documentclass{article}
\usepackage{graphicx}
\begin{document}
\title{Real-time defect detection on a KV260 platform}

\begin{abstract}
We evaluated a lightweight detector on 930 images acquired from an embedded
camera. Accuracy reached 94.2\% at 31.4573 frames per second. The obese group
showed higher latency. Detection improved by 18\%.
\end{abstract}

\section{Introduction}
As it is well known, edge inference has been related to a profile of lower
latency. Deep learning has transformed computer vision over the past decade,
with applications in medicine, autonomous driving, agriculture and industrial
inspection. However, few studies have addressed this on constrained hardware.
Therefore, there are many ways in which the pipeline can be arranged.

\section{Methods}
Images were acquired using a CMOS sensor. The obese group ($n = 930$) and the
lean group were compared. Frames were processed by the FPGA fabric. The DPU
(Deep Learning Processing Unit) was clocked at 300 MHz. We measured PSNR and
SSIM for each frame.

\section{Results}
Table~\ref{tab:acc} shows the performance of the two groups. The heavier group
achieved a mean accuracy of 94.2\% with a standard deviation of 1.2, while the
lighter group achieved a mean accuracy of 91.8\% with a standard deviation of
1.4. The mean latency of the heavier group was 12.4 $\pm$ 1.03 ms and the mean
latency of the lighter group was 15.1 $\pm$ 1.2 ms. The DPU throughput was
31.4573 frames per second.

\begin{table}[t]
\caption{Performance by group.}
\label{tab:acc}
\begin{tabular}{lcc}
group & accuracy & latency \\
\end{tabular}
\end{table}

\section{Discussion}
This study is subject to a number of limitations. The sample size was limited
and further work is required. The provided results demonstrate that SSIM values
were not inconsistent with prior work.

\end{document}
